\section{Preliminaries and Geometric Foundations}
\label{sec:preliminaries}

\subsection{Flow Matching on Euclidean Spaces}
Let $p_0$ denote a tractable base distribution on $\mathbb{R}^d$ (e.g., standard Gaussian $\mathcal{N}(0, I)$ or uniform prior) and let $p_1$ denote an empirical data distribution. Continuous Normalizing Flows (CNFs) model continuous-time transport between $p_0$ and $p_1$ through a time-dependent vector field $v_t: [0, 1] \times \mathbb{R}^d \to \mathbb{R}^d$, defining a flow $\psi_t: \mathbb{R}^d \to \mathbb{R}^d$ via the ordinary differential equation (ODE) $\frac{d}{dt} \psi_t(x) = v_t(\psi_t(x))$ with $\psi_0(x) = x$. Conditional Flow Matching (CFM) \citep{lipman2023flow, albergo2023building, liu2023flow} sidesteps the intractable marginal vector field by constructing target conditional vector fields $u_t(x_t|x_0, x_1) = x_1 - x_0$ along linear probability paths $x_t = (1 - t) x_0 + t x_1$. A neural network $v_\phi(t, x)$ is trained via the simulation-free objective:
\begin{equation}
    \mathcal{L}_{\mathrm{CFM}}(\phi) = \mathbb{E}_{t \sim \mathcal{U}[0, 1],\, x_0 \sim p_0,\, x_1 \sim p_1} \left[ \| v_\phi(t, x_t) - u_t(x_t | x_0, x_1) \|^2 \right].
\end{equation}
Once trained, new samples are generated by drawing $x_0 \sim p_0$ and integrating $v_\phi$ from $t=0$ to $t=1$ with standard numerical ODE solvers.

\subsection{The Geometry of Complex-Valued Fields: Cylinder vs. Polar Singularity}
For complex-valued physical fields, such as Magnetic Resonance Imaging (MRI) or audio spectrograms, each spatial pixel has an amplitude $A \in [0, \infty)$ and a phase $q = e^{i\theta} \in S^1$, represented by the principal angle $\theta \in [-\pi, \pi)$ when coordinates are needed; the half-open interval is a set of representatives for $S^1 \cong \mathbb{R} / 2\pi\mathbb{Z}$, not a homeomorphic copy of it. Writing $\mathbb{R}^+ = (0, \infty)$, the punctured plane is diffeomorphic to the open cylinder, $\mathbb{C}^* = \mathbb{C} \setminus \{0\} \cong \mathbb{R}^+ \times S^1$. CyFM works on its closure $[0, \infty) \times S^1$, on which the product metric below stays non-degenerate.

Standard polar coordinates on $\mathbb{C}^*$ induce the inherited Euclidean metric $g_{\mathrm{polar}} = dA^2 + A^2 d\theta^2$. This metric is flat for $A > 0$, and the parametrisation is nonsingular on its domain. What degenerates is its attempted extension to $A = 0$: the differential drops rank there, and the entire circle of phases collapses to a single point. In contrast, by equipping the single-pixel signal domain with a decoupled product metric, we define the local \textbf{cylindrical manifold}:
\begin{equation}
    \mathcal{M}_{\text{local}} = [0, \infty) \times S^1, \quad g_{\text{local}} = dA^2 + d\theta^2.
\end{equation}
Both metrics are flat (Riemann curvature tensor $R \equiv 0$), so curvature is not what separates them; $g_{\text{local}}$ differs in that the circle $A = \mathrm{const}$ has length $2\pi$ at every amplitude rather than $2\pi A$ (Appendix~\ref{sec:app_metric_choice}). The non-trivial fundamental group ($\pi_1(\mathcal{M}_{\text{local}}) = \mathbb{Z}$) is shared by both and is the price of the construction: phase must be handled modulo $2\pi$. Because the metric is decoupled, geodesics on $\mathcal{M}_{\text{local}}$ are simply the Cartesian product of geodesics on $[0, \infty)$ and minimal geodesic arcs on the circle $S^1$. Furthermore, the cylinder is parallelisable. Writing a point as $(A, q)$ with $q \in S^1 \subset \mathbb{C}$, the fields $E_A = (1, 0)$ and $E_\theta = (0, i q)$ form a global orthonormal frame, and for this metric parallel transport preserves components in it. This needs no global angular coordinate, and it removes neither the periodic wrap nor the cut locus: $\theta$ stays a local coordinate with a branch cut, which Section~\ref{sec:method} handles with a trigonometric embedding rather than with an atlas.

\subsection{Destructive Midpoint Attenuation in Euclidean Interpolation}
\label{sec:midpoint_attenuation}
When complex signals are processed as two-channel real tensors $(\mathrm{Re}(z), \mathrm{Im}(z))$, straight Euclidean paths interpolate between states $z_0 = A_0 e^{i\theta_0}$ and $z_1 = A_1 e^{i\theta_1}$. By the triangle inequality, $|z(t)| = |(1-t)z_0 + t z_1| \le (1-t) A_0 + t A_1$, with equality exactly when the two terms are non-negative multiples of one another, so the inequality is strict whenever $0 < t < 1$, both $A_0 > 0$ and $A_1 > 0$, and $\theta_0 \not\equiv \theta_1 \pmod{2\pi}$. At the temporal midpoint $t = 0.5$, the magnitude is:
\begin{equation}
    |z(0.5)| = \frac{1}{2} \sqrt{A_0^2 + A_1^2 + 2 A_0 A_1 \cos(\Delta\theta)}, \quad \Delta\theta = \theta_1 - \theta_0.
\end{equation}
For equal amplitudes $A_0 = A_1 = A$, this simplifies to $|z(0.5)| = A |\cos(\Delta\theta / 2)|$.

Write $\delta = \operatorname{wrap}(\Delta\theta) \in (-\pi, \pi]$ for the wrapped phase difference. The raw difference of two independent uniform angles is triangular on $[-2\pi, 2\pi]$, whereas its wrapped counterpart is again uniform, and $|\cos(\cdot / 2)|$ has period $2\pi$, so the midpoint magnitude depends on the endpoints through $\delta$ alone.

\begin{proposition}[Analytic Midpoint Attenuation]
\label{prop:attenuation}
For equal amplitudes $A_0 = A_1 = A > 0$ and a uniform wrapped phase difference $\delta \sim \mathcal{U}(-\pi, \pi]$, which an independent coupling of uniform phases induces, the relative expected midpoint amplitude under Euclidean interpolation suffers an analytic attenuation of $1 - \frac{2}{\pi} \approx 36.338\%$:
\begin{equation}
    \frac{\mathbb{E}[|z(0.5)|]}{A} = \frac{1}{2\pi} \int_{-\pi}^{\pi} \left| \cos\left(\frac{\delta}{2}\right) \right| d\delta = \frac{2}{\pi} \approx 0.6366.
\end{equation}
\end{proposition}


For $A_0, A_1 > 0$, exactly antipodal endpoints are the only ones whose path meets the origin, and then only at the single time $t = A_0 / (A_0 + A_1)$, where the phase is undefined; under any continuous phase law that has probability zero. A vanishing endpoint gives a radial chord instead, whose argument never turns. The practical case is the near-antipodal one at comparable amplitudes, where the path passes arbitrarily close to zero: the amplitude is heavily attenuated and the induced angular velocity is correspondingly large. The attenuation itself occurs for any $\delta \neq 0$; its expected value $2/\pi$ is the figure for an independent coupling, and it moves with the coupling, since minibatch OT suppresses large $|\delta|$ (Table~\ref{tab:bridge_ablation}). That dependence is itself the asymmetry between the two geometries. What a Cartesian path costs can only be stated for a given law of $\delta$, and has to be bought down by choosing a coupling; the cylindrical bound $|u_\theta| \le \pi$ holds for every pair of endpoints under every coupling, with no distributional assumption at all.